%!TEX root = ../report.tex
\documentclass[report.tex]{subfiles}
\begin{document}
    \begin{abstract}
        
        Reliable state estimation is vital for \gls{rwr} operating in unstructured planetary environments where external positioning systems such as \gls{gnss} and motion capture systems are unavailable. One of the main challenges in dead-reckoning odometry approaches is the significant error accumulation due to sensor drift and unaccounted wheel slippage on loose terrains such as sand and gravel. This thesis presents an odometry approach using sensor fusion techniques within an \gls{inekf} framework. This approach is developed for the \gls{rwr} platform, where the unique \gls{rw} geometry and the associated wheel-ground contact dynamics make the existing baseline odometry susceptible to wheel slippage and drift. The measured \gls{pcv} is calculated using the joint encoder measurements and rimless-wheel spoke geometry, while the predicted \gls{pcv} is calculated using the current filter state and the angular velocity measured by the \gls{imu}. These quantities are then used in the update step of the \gls{inekf} to estimate the rover state. The inconsistent joint encoder measurements compared to the \gls{imu} measurements are rejected before the update step of the \gls{inekf}. 
        
        This prevents the erroneous joint encoder measurements from being incorporated into the state estimation process. The pose estimates 
        from the slip-aware odometry estimator are evaluated against the pose estimates from the motion capture system which is considered as the ground truth. The proposed approach is also compared with the baseline odometry currently integrated into the Coyote~3 \gls{rwr} platform. The evaluation is performed through extensive real-world experiments with different trajectories on indoor floor and sand terrain. The evaluation results are presented as the translational \gls{rpe} drift \gls{rmse}, rotational \gls{rpe} drift \gls{rmse}, and other statistical performance metrics. 
        
        The results show that rejecting inconsistent \gls{pcv} measurements when slip is detected improves the odometry estimation in terms of translational \gls{rpe} drift \gls{rmse} for the crater-terrain dataset. On the sand-terrain datasets, the proposed approach achieves a lower translational \gls{rpe} drift \gls{rmse} in six of the seven selected datasets. However, on the flat-floor datasets, the baseline odometry performs better than the proposed approach in terms of \gls{rpe} metrics. 
        
    \end{abstract}
\end{document}
